\chapter{Estudio de extremos locales. Condiciones suficientes}

Dado $\cc{F}:D \to \bb{R}$, dado $y_*\in D$ un punto crítico, estudiar su carácter de máximo o mínimo local se reduce a estudiar
\begin{gather*}
    \veps \mapsto \cc{F}[y_* + \veps \varphi] \ \ \text{ para cualquier $\varphi$ admisible}
\end{gather*}
De hecho, 
\begin{gather*}
    \cc{F}[y_*+\veps\varphi] = \cc{F}[y_*] + \frac{\veps^2}{2} \frac{d^2}{d\veps^2} \cc{F}[y_* + \veps\varphi]_{|_{\veps = 0}} + O(\veps^3)
\end{gather*}
Solamente consideramos el siguiente caso:\\

$I=[x_0, x_1]$, $F\in C^3(I \times \bb{R} \times \bb{R})$ dada por 
\begin{gather*}
    \cc{F}[y] = \int_{x_0}^{x_1} F(x, y(x), y'(x)) dx
\end{gather*}
definido sobre
\begin{gather*}
    D = \{y\in C^1(I) : y(x_0)=y_0,\ y(x_1=y_1)\}
\end{gather*}
con $y_0,y_1\in \bb{R}$ dados. Las direcciones admoisibles son $\varphi\in C_0^1(I)$.\\

Dado $y\in D$, $\varphi$ admisible, formamos
\begin{gather*}
    y_\veps := y +\veps \varphi
\end{gather*}
Una vez hecho eso calculamos la derivada segunda
\begin{gather*}
    \frac{d^2}{d\veps^2}\cc{F}[y_\veps] = \frac{d}{d\veps} \int_{x_0}^{x_1} \frac{d}{d\veps} [F(x, y_\veps(x), y_\veps'(x))] dx
\end{gather*}
Utilizaremos la notación que venimos usando en el tema anterior:
\begin{gather*}
    F=F(x,y,p),\ \ \ \ F_y=\frac{\partial F}{\partial y},\ \ \ \ F_p = \frac{\partial F}{\partial p}
\end{gather*}
y la expresión anterior quedará como 
\begin{align*}
    \frac{d^2}{d\veps^2}\cc{F}[y_\veps] &= \frac{d}{d\veps} \int_{x_0}^{x_1} F_y(x,y_\veps(x), y_\veps'(x))\varphi(x) + F_p(x,y_\veps(x), y_\veps'(x))\varphi'(x)dx = \\
    &= \int_{x_0}^{x_1} F_{yy}(x,y_\veps(x), y_\veps'(x)) \varphi^2(x) + F_{yp}(x,y_\veps(x), y_\veps'(x))\varphi'(x)\varphi(x) + \\
    &+ F_{py}(x,y_\veps(x), y_\veps'(x))\varphi(x)\varphi'(x) + F_{pp}(x,y_\veps(x), y_\veps'(x))\varphi'(x)^2 dx
\end{align*}
Como $F\in C^3(I\times \bb{R} \times \bb{R})$ tendremos que $F_{yp}=F_{py}$. Además
\begin{gather*}
    2\int_{x_0}^{x_1} F_{yp}(x,y_\veps(x), y_\veps'(x)) \varphi(x)\varphi'(x) dx = \int_{x_0}^{x_1} F_{yp}(x,y_\veps(x), y_\veps'(x)) \frac{d}{dx}(\varphi(x))^2 dx
\end{gather*}
Si integramos ahora por partes tenemos que 
\begin{gather*}
    \int_{x_0}^{x_1} F_{yp}(x,y_\veps(x), y_\veps'(x)) \frac{d}{dx}(\varphi(x))^2 dx = -\int_{x_0}^{x_1} \frac{d}{dx} [F_{yp}(x,y_\veps(x), y_\veps'(x))] \varphi(x)^2 dx
\end{gather*}
Si ahora unimos esto a la expresión anterior llegamos a 
\begin{align*}
    \frac{d^2}{d\veps^2}\cc{F}[y_\veps] &= \int_{x_0}^{x_1} F_{pp}(x,y_\veps(x), y_\veps'(x)) \varphi'(x)^2 +\\
    &+ \left(-\frac{d}{dx}[F_{yp}(x,y_\veps(x), y_\veps'(x)) ] + F_{yy}(x,y_\veps(x), y_\veps'(x))\right)\varphi(x)^2 dx
\end{align*}
Por tanto, al evaluar en $\veps = 0$ tenemos
\begin{gather*}
    \frac{d^2}{d\veps^2}\cc{F}[y_\veps]_{|_{\veps=0}} =\\
    \int_{x_0}^{x_1} F_{pp}(x,y(x), y'(x)) \varphi'(x)^2 + \left(-\frac{d}{dx}[F_{yp}(x,y(x), y'(x)) ] + F_{yy}(x,y(x), y'(x))\right)\varphi(x)^2 dx
\end{gather*}
En varias fuentes podremos encontrar la siguiente noctación:
\begin{gather*}
    P = F_{pp}(x,y(x), y'(x)),\ \ \ \ Q = -\frac{d}{dx}[F_{yp}(x,y(x), y'(x)) ] + F_{yy}(x,y(x), y'(x))
\end{gather*}

\begin{definicion}[Condición necesaria de Legendre]
    Diremos que el punto crítico $y_*\in D$ satisface la \textbf{condición necesaria de Legendre} para minimalidad si 
    \begin{gather*}
        P \geq 0\ \ \ \forall x \in  I
    \end{gather*}
    Si esto no se cumple, entonces $y_*$ no será un mínimo\footnote{no veremos la prueba pero se puede encontrar en la bibliografía}. Diremos que es estricta si 
    \begin{gather*}
        P > 0\ \ \ \forall x \in  I
    \end{gather*}
\end{definicion}

De forma sintética escribiremos
\begin{gather*}
    \frac{d^2}{d\veps^2}\cc{F}[y_\veps]_{|_{\veps=0}} = \int_{x_0}^{x_1} P(\varphi')^2 + Q\varphi^2 dx
\end{gather*}
Observamos que 
\begin{gather*}
    \int_{x_0}^{x_1} (P\varphi')\varphi' dx = -\int_{x_0}^{x_1} \frac{d}{dx} (P\varphi'(x)) \varphi(x) dx
\end{gather*}
Por tanto,
\begin{gather*}
    \frac{d^2}{d\veps^2}\cc{F}[y_\veps]_{|_{\veps=0}} = \int_{x_0}^{x_1} \left( -\frac{d}{dx} (P\varphi') + Q\varphi \right) \varphi dx
\end{gather*}

\begin{definicion}[Ecuación de Jacobi]
    Dado $y\in D$ punto crítico, diremos que su \textbf{ecuación de Jacobi asociada} es 
    \begin{gather*}
        -\frac{d}{dx}(P\varphi') + Q\varphi = 0 \ \ \text{ en } [x_0, x_1]
    \end{gather*}
    La incógnita en esta ecuación es $\varphi\in C^2([x_0, x_1])$.
\end{definicion}

\begin{definicion}[Punto conjugado]
    Diremos que un cierto punto $\bar{x}\in ]x_0, x_1]$ es un \textbf{punto conjugado} al $x_0$ si podemos encontrar una solución $\varphi$ de la ecuación de Jacobi tal que  se verifiquen las siguientes condiciones
    \begin{itemize}
        \item $\varphi \not\equiv 0$
        \item $\varphi(x_0)=0$
        \item $\varphi(\bar{x})=0$
    \end{itemize}
\end{definicion}

\begin{teo}[Condición de Jacobi]
    Sea $I=[x_0,x_1]$, $F\in C^3(I \times \bb{R} \times \bb{R})$. Consideramos 
    \begin{gather*}
        \cc{F}[y] = \int_{x_0}^{x_1} F(x,y(x),y'(x))dx
    \end{gather*}
    definido sobre $D = \{y\in C^1(I) : y(x_0)=y_0,\ y(x_1=y_1)\}$ con $y_0, y_1\in \bb{R}$ dados. Sea $y\in D$ un punto crítico que verifique la condición de Legendre estricta, es decir, $P>0$ para todo $x\in I$. Entonces:
    \begin{itemize}
        \item (Condición necesaria) Si $y$ es mínimo local estricto, entonces no existen puntos conjugados a $x_0$ en el intervalo $]x_0, x_1[$.
        \item (Condición suficiente) Si $x_0$ no tiene puntos conjugados en $]x_0,x_1]$, entonces $y$ es un mínimo local estricto.
    \end{itemize}
\end{teo}

\begin{coro}
    Bajo estas circunstancias, si $[x_0,x_1]$ es suficientemente pequeño\footnote{dependerá de cada problema concreto.}, se cumple la condición de Jacobi.
\end{coro}

\begin{ejemplo}
    Es espacio ambiente es $X=C^1(I)$ con la norma $\|u\|_{C^1} = \|u\|_\infty + \|u'\|_\infty$. Estudiamos $\cc{F}[y] = \int_0^1 (y'(x))^3dx$ sobre
    \begin{gather*}
        D=\{y\in C^1([0,1]) : y(0)=0,\ y(1)=1\}
    \end{gather*}
    Notaremos $x_0=0$ y $x_1=1$ (para adaptar el enunciado a nuestra notación). Como el funcional solo depende de $p$ tendremos $F=F(p)=p^3$. La ecuación de Euler-Lagrange queda
    \begin{gather*}
        -\frac{d}{dx}(F_p) = 0 \Rightarrow \frac{d}{dx} [3(y'(x))^2] = 0
    \end{gather*}
    lo que nos dice que $y'(x)^2 = cte$, luego $y'(x)=cte$. Por tanto, como $y\in D$ tendremos que los extremales son de la forma
    \begin{gather*}
        y(x)=c_1 + c_2x,\ \ c_1,c_2\in \bb{R}
    \end{gather*}
    De la condición $y\in D$ llegamos a que la única extremal es $y(x)=x$. Ahora nos podríamos preguntar si esto es un máximo o un mínimo local. Tenemos que 
    \begin{gather*}
        F_{pp} = 6p
    \end{gather*}
    Para $y(x)=x$ tenemos $y'(x)=1$ y por tanto $P=6>0$ para todo $x\in [x_0, x_1]$ (se cumple la condición de Legendre estricta). Estamos en condiciones de comprobar la condición de Jacobi. La ecuación de Jacobi que es
    \begin{gather*}
        -\frac{d}{dx}(P\varphi') + Q\varphi = 0
    \end{gather*}
    Como en este caso $Q=0$ tendremos que la ecuación de Jacobi será
    \begin{gather*}
        -6\varphi''=0
    \end{gather*}
    Buscamos soluciones de la forma $\varphi(x_0)=0$ y $\varphi(\bar{x})=0$ para algún $x_0<\bar{x}\leq x_1$. De nuevo, la solución general es $\varphi(x)=c_1+c_2x$ con $c_1,c_2\in \bb{R}$. Al imponer $\varphi(x_0)=0$ llegamos a que $c_1=0$, de forma que $\varphi(x)=c_2x$ para $c_2\in \bb{R}$. Nos preguntamos ahora si existe algún $\bar{x}>0$ tal que $\varphi(\bar{x})=0$. La respuesta es que independientemente del valor de $c_2\in \bb{R}$ no existe un tal $\bar{x}$ ($\varphi$ es lineal). Por tanto, no hay puntos conjugados, lo que nos dice que la condición de Jacobi se cumple. Por tanto $y(x)=x$ es un mínimo local estricto.
\end{ejemplo}

\begin{ejemplo}
    Estudiamos $\cc{F}[y] = \int_0^a (y')^2 - y^2 dx$ en $D=C_0^1([0,a])$.\\

    Al aplicar la notación estudiada tenemos 
    \begin{gather*}
        F=p^2 - y^2,\ \ \ \ F_p=2p,\ \ \ \ F_{y}=-2y,\ \ \ \ F_{pp}=2,\ \ \ \ F_{yy}=-2,\ \ \ \ F_{py}=0
    \end{gather*}
    La ecuación de Euler-Lagrange es en este caso 
    \begin{gather*}
        -2y(x) - \frac{d}{dx}(2y'(x)) = 0
    \end{gather*}
    Lo que nos lleva a la ecuación diferencial
    \begin{gather*}
        y''(x)+y(x)=0
    \end{gather*}
    Por lo que las extremales serán de la forma 
    \begin{gather*}
        y(x) = A\sin(x) + B\cos(x),\ \ \ A,B\in \bb{R}
    \end{gather*}
    Como $y(0)=0$ llegamos a que $B=0$. Por tanto $y(x)= A\sin(x)$. O bien $A=0$ o bien $a=u\pi$ para $u\in \bb{Z}$. Es decir, en general el único punto crítico es $y(x)=0$. Para cualquier punto $y\in D$ tenemos que $P=2>0$ por lo que se cumple la ecuación de Legendre estricta (para mínimo local). Nos preguntamos ahora si hay puntos conjugados para $x_0=0$, cuestión que nos lleva a calcular la Ecuación de Jacobi
    \begin{gather*}
        -2\varphi''-2\varphi = 0
    \end{gather*}
    La solución general es $\varphi(x) = C\cos(x) + D\sin(x)$ con $C,D\in \bb{R}$. Si $\varphi(0)=0$, entonces $D=0$, por lo que $\varphi(x) = C\cos(x)$. Buscamos ahora la posibilidad de que exista un $0<\bar{x}\leq a$ tal que $\varphi(\bar{x})=0$, entonces o bien $C=0$, lo cual no aporta información, o bien $\bar{x}=n\pi$ con $n\in \bb{N}$. Si $a\geq \pi$, entonces podemos encontrar algún $\bar{x}\leq a$ en esas circunstancias. Cabe hacer una distinción de casos:
    \begin{itemize}
        \item Si $a<\pi$, no hay puntos conjugados por lo que $y(x)=0$ es un mínimo local estricto.
        \item Si $a>\pi$, la función $y(x)=0$ no es un mínimo local estricto (se viola la condición de Jacobi) 
    \end{itemize}
    En el caso de que se dé la igualdad es más complicado extraer una conclusión.
\end{ejemplo}

\section{Funcionales integrales definidos sobre curvas}
Dado $I\subset \bb{R}$ un intervalo, digamos $I=[x_0, x_1]$, nos interesan las curvas parametrizadas $y:I \to \bb{R}^N$. Entendemos que $y=(y_1,y_2,...,y_n)$, con $y_i:I\to \bb{R}$ para cada $1\leq i \leq N$. Definimos 
\begin{gather*}
    C(I;\bb{R}^N):= \{y:I\to \bb{R}^N : y\in C(I,\bb{R})\text{ para } 1\leq i \leq N\}\\
    C^1(I;\bb{R}^N):= \{y:I\to \bb{R}^N : y\in C^1(I,\bb{R})\text{ para } 1\leq i \leq N\}\\
    C_0^1(I;\bb{R}^N):= \{y:I\to \bb{R}^N : y\in C^1(I,\bb{R})\text{ para } 1\leq i \leq N,\ y(x_0) = \vec{0},\ y(x_1) = \vec{0}\}
\end{gather*}
Si el contexto lo deja claro podremos escribir $C(I)$, $C^1(I)$, etc.\\

Dada $y:I \to \bb{R}^N$ nos interesarán funcionales de la forma
\begin{gather*}
    \cc{F} = \int_{x_0}^{x_1} F(x, y(x), y'(x)) dx
\end{gather*}
Se entiende que $y':I \to \bb{R}^N$ viene dado por 
\begin{gather*}
    y'(x) = (y_1'(x), y_2'(x), ..., y_N'(x)) \ \ \ \ \forall x \in I
\end{gather*}
y además por su definición se entiende que $F:I\times \bb{R}^N \times \bb{R}^N \to \bb{R}$.\\

\begin{notacion}
    Digamos que $F=F(x,y,p)$ donde $y=(y_1, y_2,...,y_N)$ y $p = (p_1,p_2,...,p_N)$ donde $p_i = y_i'$ para cada $1\leq i \leq N$. Introduciremos de nuevo los atajos de notación que introducíamos para el caso escalar
    \begin{gather*}
        F_y := \left(\frac{\partial F}{\partial y_1}, \frac{\partial F}{\partial y_2}, ..., \frac{\partial F}{\partial y_N}\right),\ \ \ \
        F_p := \left(\frac{\partial F}{\partial p_1}, \frac{\partial F}{\partial p_2}, ..., \frac{\partial F}{\partial p_N}\right)
    \end{gather*}
\end{notacion}

El caso que trabajaremos típicamente será un $F\in C^2(I\times \bb{R}^N \times \bb{R}^N)$ dado por un funcional $\cc{F}$ definido sobre
\begin{gather*}
    D = \{y\in C^1(I;\bb{R}^N) : y(x_0) = y_0,\ y(x_1)=y_1\}
\end{gather*}
con $y_0, y_1\in \bb{R}^N$ dados.\\

Dado $y_*\in C^2(I;\bb{R}^N)\cap D$, si es punto crítico verifica la ecuación de Euler-Lagrange, es decir, 
\begin{gather*}
    F_y(x,y_*(x), y_*'(x)) - \frac{d}{dx} [F_p(x,y_*(x), y_*'(x))] = \vec{0}\ \ \ \ \forall x \in I
\end{gather*}
que será un sistema de $N$ ecuaciones diferenciales ordinarias. Al calcular $\delta_y\cc{F}[y]$ se obtiene 
 \begin{gather*}
    \delta_y\cc{F}[y] = \int_{x_0}^{x_1} F_y(x,y(x), y'(x)) \varphi(x) + F_p(x,y(x, y'(x))) \varphi'(x) dx
 \end{gather*}
 Las direcciones admisibles son $\varphi \in C_0^1(I:\bb{R}^N)$ donde seguimos definiendo $y_\veps := y + \veps \varphi$.
 \begin{align*}
    \int_{x_0}^{x_1} \frac{d}{d\veps}(F(x,y_\veps, y'_\veps)) dx &= \int_{x_0}^{x_1} \frac{\partial F(x,y_\veps, y'_\veps)}{y_1} \frac{d(y_\veps)_1}{d \veps} + \frac{\partial F(x,y_\veps, y'_\veps)}{y_2} \frac{d(y_\veps)_2}{d \veps} + ... = \\
    &= \int_{x_0}^{x_1} \frac{\partial F(x,y_\veps, y'_\veps)}{y_1} \varphi_1 + \frac{\partial F(x,y_\veps, y'_\veps)}{y_2} \varphi_2 + ... = \\
    &= \int_{x_0}^{x_1} F_y(x,y(x), y'(x)) \varphi(x)
 \end{align*}
 Queremos ahora integrar por partes en el término
 \begin{gather*}
    F_p(x,y(x, y'(x))) \varphi'(x) = \sum_{i=1}^{N}F_{p_i}(x, y(x), y'(x))\varphi'_i
 \end{gather*}
 Véase que 
 \begin{gather*}
    \int_{x_0}^{x_1} F_{p_1}(x, y(x), y'(x)) \varphi'_1(x) dx = \\
    - \int_{x-0}^{x_1} \frac{d}{dx} [F_{p_1}(x,y(x), y'(x))]\varphi_1(x) dx + \left[ F_{p_1}(x,y(x), y'(x)) \varphi_1(x)\right]_{x=x_0}^{x=x_1} 
 \end{gather*}
 Recordemos que $\varphi(x_0) = \vec{0} = \varphi(x_1)$ por lo que 
 \begin{gather*}
    \int_{x_0}^{x_1} F_{p_1}(x, y(x), y'(x)) \varphi'_1(x) dx =
    - \int_{x-0}^{x_1} \frac{d}{dx} [F_{p_1}(x,y(x), y'(x))]\varphi_1(x) dx
 \end{gather*}
 y por tanto 
 \begin{align*}
    \int_{x_0}^{x_1} F_p(x, y(x), y'(x)) \varphi'(x) dx &= - \int_{x_0}^{x_1} \sum_{i=1}^N \frac{d}{dx} [F_{p_i}(x, y(x), y'(x))] \varphi_i(x) dx =\\
    &= -\int_{x_0}^{x_1} \frac{d}{dx} [F_p(x, y(x), y'(x))]\varphi(x) dx
 \end{align*}

 Al unir todos estos resultados nos quedará que si $y\in I$ es punto crítico, entonces
 \begin{gather*}
    \int_{x_0}^{x_1} [F_y(x, y(x), y'(x)) - \frac{d}{dx}F_p(x, y(x), y'(x))] \varphi(x)dx = 0\ \ \ \ \forall \varphi\in C_0^1(I;\bb{R}^N)
 \end{gather*}

 \begin{lema}[Lema fundamental. Versión para curvas]
    Dada $f\in C([x_0, x_1]; \bb{R}^N)$, si 
    \begin{gather*}
        \int_{x_0}^{x_1} f \cdot \varphi \ dx = 0 \ \ \ \ \forall \varphi \in C_c^\infty(]x_0, x_1[; \bb{R}^N)
    \end{gather*}
    entonces $f=\vec{0}$ para todo $x\in [x_0, x_1]$.
    \begin{proof}
        Tenemos que 
        \begin{gather*}
            f\cdot \varphi = f_1 \cdot \varphi_1 + f_2 \cdot \varphi_2 + ... + f_N\cdot \varphi_N
        \end{gather*}
        Escojo $\varphi = (\varphi_1, 0, 0, ..., 0) \in C_c^\infty(]x_0, x_1[;\bb{R}^N)$ y deducimos que $f_1(x) = 0$ para todo $x\in I$. Si repetimos este razonamiento de forma análoga para los restantes $f_i$ obtenemos el resultado del lema.
    \end{proof}
 \end{lema}

